% Durée indicative de la section : 3 min 30 à 4 min 30.

% ---------- Périmètre de la base (~45 s) ----------
\begin{frame}{Une base de 14\,649 séjours opératoires}
  \begin{center}
    \small Une ligne correspond à un \textbf{cas hospitalier associé à une intervention}.
  \end{center}

  \begin{columns}[T]
    \begin{column}{0.31\textwidth}
      \begin{exampleblock}{\centering Volume}
        \centering
        {\LARGE\textbf{14\,649}}\\[-1pt]
        {\small cas hospitaliers}
      \end{exampleblock}
    \end{column}
    \begin{column}{0.31\textwidth}
      \begin{block}{\centering Patients}
        \centering
        {\LARGE\textbf{11\,108}}\\[-1pt]
        {\small identifiants distincts}
      \end{block}
    \end{column}
    \begin{column}{0.31\textwidth}
      \begin{alertblock}{\centering Variables}
        \centering
        {\LARGE\textbf{30}}\\[-1pt]
        {\small administratives, cliniques et bloc}
      \end{alertblock}
    \end{column}
  \end{columns}

  \vspace{1em}
  \begin{columns}[T]
    \begin{column}{0.48\textwidth}
      \footnotesize
      \textbf{Administratif et clinique}\par
      Dates (2019--2022), âge, sexe, identifiant patient,
      diagnostics CIM-10 et actes CCAM.
    \end{column}
    \begin{column}{0.48\textwidth}
      \footnotesize
      \textbf{Parcours opératoire}\par
      Praticien, chirurgien, type d'intervention, anesthésie
      et horaires clés du bloc.
    \end{column}
  \end{columns}

  \vspace{0.6em}
  {\footnotesize\textcolor{mcmAlert}{\textbf{Point de vigilance :}}
  la granularité est le \textbf{séjour}, pas le patient ;
  2\,546 patients apparaissent dans plusieurs cas.\par}
\end{frame}

% ---------- Qualité des données (~70 s) ----------
\begin{frame}{Qualité : deux anomalies structurantes}
  \begin{columns}[T]
    \begin{column}{0.48\textwidth}
      \begin{alertblock}{Durée de séjour}
        \small
        \begin{itemize}\itemsep3pt
          \item \textbf{2\,560 incohérences}, soit \textbf{17,5\,\%}
          \item 1\,884 séjours notés « 1 jour » malgré au moins une nuit
          \item 31 valeurs proches de dates juliennes
        \end{itemize}
      \end{alertblock}

      \begin{block}{Référence fiable}
        \footnotesize Aucune date manquante ; la sortie ne précède jamais
        l'entrée et l'intervention est toujours comprise dans le séjour.
      \end{block}
    \end{column}

    \begin{column}{0.48\textwidth}
      \begin{exampleblock}{Horaires opératoires à zéro}
        \centering\scriptsize
        \begin{tabular}{@{}lr@{}}
          SAS pré-anesthésie & 133 \\
          Entrée en salle & 172 \\
          Incision & 552 \\
          Sortie de salle & 172 \\
        \end{tabular}
      \end{exampleblock}

      \begin{block}{Inversions chronologiques}
        \small
        \begin{itemize}\itemsep3pt
          \item 42 incisions avant l'entrée en salle
          \item 4 sorties avant l'incision
          \item aucune sortie avant l'entrée en salle
        \end{itemize}
      \end{block}
    \end{column}
  \end{columns}

  \begin{center}
    \small\textcolor{mcmPrimary}{\textbf{Ces inversions sont des erreurs de saisie,
    pas des opérations franchissant minuit.}}
  \end{center}
\end{frame}

% ---------- Nettoyage retenu (~65 s) ----------
\begin{frame}[t]{Nettoyage : 14\,507 lignes exploitables}
  \begin{center}
    \begin{tikzpicture}[
      node distance=0.55cm,
      every node/.style={font=\small, align=center},
      box/.style={draw=mcmPrimary, very thick, rounded corners=2pt,
                  minimum height=1.25cm, text width=3.0cm},
      processbox/.style={draw=mcmExample, very thick, rounded corners=2pt,
                        fill=mcmExample!12, minimum height=1.25cm, text width=3.7cm}
    ]
      \node[box] (raw) {\Large\textbf{14\,649}\\\footnotesize lignes sources};
      \node[processbox, right=of raw] (filter) {\textbf{Filtre qualité}\\[-1pt]
        \footnotesize type d'intervention ou praticien absent};
      \node[box, right=of filter, fill=mcmPrimaryMuted] (clean) {
        \Large\textbf{14\,507}\\\footnotesize lignes conservées};
      \draw[-{Latex}, very thick, mcmPrimary] (raw) -- (filter);
      \draw[-{Latex}, very thick, mcmPrimary] (filter) -- (clean);
    \end{tikzpicture}
  \end{center}

  \begin{columns}[T]
    \begin{column}{0.47\textwidth}
      \begin{block}{Suppression ciblée}
        \small
        \begin{itemize}\itemsep2pt
          \item 137 types d'intervention absents
          \item 98 praticiens absents
          \item 142 lignes ont au moins 1 des 2 absences
        \end{itemize}
      \end{block}
    \end{column}
    \begin{column}{0.49\textwidth}
      \begin{exampleblock}{Durée corrigée}
        \small
        \[
          d_{\mathrm{corrigée}}=
          (\mathrm{sortie}-\mathrm{entrée})+1
        \]
        \footnotesize 2\,488 durées corrigées parmi les lignes conservées.
      \end{exampleblock}
    \end{column}
  \end{columns}

  \vspace{0.35em}
  {\footnotesize\textcolor{mcmAlert}{\textbf{Traçabilité :}}
  la durée source est conservée ; une colonne corrigée et un indicateur
  d'incohérence sont ajoutés ; la table nettoyée est exportée séparément.\par}
\end{frame}

% ---------- Exploitabilité et limites (~60 s) ----------
\begin{frame}{Ce que la base permet --- et ne permet pas}
  \begin{columns}[T]
    \begin{column}{0.48\textwidth}
      \begin{exampleblock}{Disponible pour prédire}
        \small
        \begin{itemize}\itemsep3pt
          \item durée de séjour corrigée ;
          \item temps d'occupation de salle après validation ;
          \item âge, sexe, diagnostics et actes ;
          \item praticien et type d'intervention.
        \end{itemize}
      \end{exampleblock}
    \end{column}
    \begin{column}{0.48\textwidth}
      \begin{alertblock}{Absent pour ordonnancer}
        \small
        \begin{itemize}\itemsep3pt
          \item identifiant des salles et vacations ;
          \item capacité journalière en lits ;
          \item disponibilité des équipes ;
          \item urgences, annulations et refus de dates.
        \end{itemize}
      \end{alertblock}
    \end{column}
  \end{columns}

  \vspace{0.25em}
  \begin{block}{Conclusion de l'analyse}
    \centering
    La base nettoyée suffit pour \textbf{apprendre les durées}.\\
    L'optimisation nécessitera des \textbf{données métier complémentaires}
    ou des hypothèses de simulation explicites.
  \end{block}
\end{frame}
